\subsection{SEMI-DIRECT PRODUCTS}

Let L, N be two subgroups of a group G, such that LN = NL; then LN is a subgroup of G, since

\[
(\mathrm{LN}) (\mathrm{LN}) ^ {- 1} = \mathrm{LNN} ^ {- 1} \mathrm{L} ^ {- 1} = \mathrm{LNL} = \mathrm{LLN} = \mathrm{LN}
\]

Moreover, the mapping \(\varphi: (x, y) \to xy\) of \(N \times L\) into \(G\) is injective if and only if \(N \cap L = \{e\}\) . For it is clear that if \(\varphi\) is injective, then \(N \cap L = \{e\}\) ; and conversely, the relation \(x' y' = xy\) , where \(x, x' \in N\) and \(y, y' \in L\) , implies that \(x^{-1}x' = yy'^{-1} \in N \cap L\) ; hence if \(N \cap L = \{e\}\) , then \(\varphi\) is injective. Thus \(\varphi\) is bijective if and only if \(LN = G\) and \(N \cap L = \{e\}\) .

If N is a normal subgroup of G (or more generally is normal in some subgroup of G containing \(N \cup L\) ), the condition LN = NL is automatically satisfied. Moreover, for each \(y \in L\) the mapping \(\sigma_{y}: x \to yxy^{-1}\) is an automorphism of the group N, and for any two elements u, v of L we have

\[
\sigma_ {u v} = \sigma_ {u} \circ \sigma_ {v};\tag{1}
\]

and for any \(x, x'\) in \(\mathbf{N}\) and \(y, y'\) in \(\mathbf{L}\) we have also

\[
(x y) (x ^ {\prime} y ^ {\prime}) = (x \sigma_ {y} (x ^ {\prime})) (y y ^ {\prime}).\tag{2}
\]

Conversely :

PROPOSITION 27. Let N and L be two groups, \(e'\) and \(e''\) their respective identity elements. Suppose that we are given a homomorphism \(y \to \sigma_y\) of L into the group \(\Gamma\) of automorphisms of N. Then:

\begin{enumerate}
\def\labelenumi{\arabic{enumi})}
\tightlist
\item
  On the product set \(\mathbf{S} = \mathbf{N} \times \mathbf{L}\) , the internal law of composition
\end{enumerate}

\[
(x, y) (x ^ {\prime}, y ^ {\prime}) = (x \sigma_ {y} (x ^ {\prime}), y y ^ {\prime})\tag{3}
\]

defines a group structure, for which \(j_1: x \to (x, e'')\) is an isomorphism of \(N\) onto a normal subgroup of \(S\) , \(j_2: y \to (e', y)\) is an isomorphism of \(L\) onto a subgroup of \(S\) , and \(\mathrm{pr}_2: S \to L\) is a surjective homomorphism whose kernel is \(j_1(N)\) and which is such that \(\mathrm{pr}_2 \circ j_2\) is the identity automorphism of \(L\) .

\begin{enumerate}
\def\labelenumi{\arabic{enumi})}
\setcounter{enumi}{1}
\tightlist
\item
  Let \(f: \mathbf{N} \to \mathbf{G}\) , \(g: \mathbf{L} \to \mathbf{G}\) be two homomorphisms into a group \(\mathbf{G}\) , such that
\end{enumerate}

\[
f (\sigma_ {y} (x)) = g (y) f (x) g (y ^ {- 1})\tag{4}
\]

for all \(x \in \mathbb{N}\) and all \(y \in \mathbb{L}\) . Then there is a unique homomorphism \(h: \mathbb{S} \to \mathbb{G}\) such that \(f = h \circ j_1\) , and \(g = h \circ j_2\) .

If \(x, x', x''\) are elements of N and \(y, y', y''\) are elements of L, then

and

\[
\begin{array}{r l} ((x, y) (x ^ {\prime}, y ^ {\prime})) (x ^ {\prime \prime}, y ^ {\prime \prime}) & = (x \sigma_ {y} (x ^ {\prime}), y y ^ {\prime}) (x ^ {\prime \prime}, y ^ {\prime \prime}) \\ & = (x \sigma_ {y} (x ^ {\prime}) \sigma_ {y y ^ {\prime}} (x ^ {\prime \prime}), y y ^ {\prime} y ^ {\prime \prime}) \\ (x, y) ((x ^ {\prime}, y ^ {\prime}) (x ^ {\prime \prime}, y ^ {\prime \prime})) & = (x, y) (x ^ {\prime} \sigma_ {y ^ {\prime}} (x ^ {\prime \prime}), y ^ {\prime} y ^ {\prime \prime}) \\ & = (x \sigma_ {y} (x ^ {\prime} \sigma_ {y ^ {\prime}} (x ^ {\prime \prime})), y y ^ {\prime} y ^ {\prime \prime}) \end{array}
\]

and therefore the associativity of the law (3) follows from the facts that \(y \to \sigma_{y}\) is a homomorphism of L into \(\Gamma\) and that \(\sigma_{y}\) is an automorphism of N. Clearly \((e', e'')\) is the identity element of (3), and finally

\[
(x, y) \left(\sigma_ {y ^ {- 1}} (x ^ {- 1}), y ^ {- 1}\right) = \left(\sigma_ {y ^ {- 1}} (x ^ {- 1}), y ^ {- 1}\right) (x, y) = \left(e ^ {\prime}, e ^ {\prime \prime}\right)
\]

so that \((x, y)\) has an inverse in S. The other assertions of 1) are clear. On the other hand, since \((x, y) = (x, e'') (e', y)\) , a homomorphism \(h\) which satisfies the conditions of 2) must necessarily satisfy

\[
h (x, y) = f (x) g (y),
\]

hence is unique if it exists; moreover it is immediate from (4) that

\[
f \left(x \sigma_ {y} (x ^ {\prime})\right) g (y y ^ {\prime}) = f (x) g (y) f (x ^ {\prime}) g (y ^ {- 1}) g (y) g (y ^ {\prime}) = f (x) g (y) f (x ^ {\prime}) g (y ^ {\prime})
\]

which shows that \((x, y) \to f(x)g(y)\) is indeed a homomorphism of S into G satisfying the conditions of 2).

COROLLARY. The homomorphism \(h\) defined in 2) of Proposition 27 is injective if and only if \(f\) and \(g\) are injective and \(f(\mathbf{N}) \cap g(\mathbf{L}) = \{e\}\) ; and \(h\) is surjective if and only if \(f(\mathbf{N})g(\mathbf{L}) = \mathbf{G}\) .

Since \(h(x, y) = f(x)g(y)\) , the second assertion is obvious; moreover it follows from (4) that \(f(\mathbf{N})g(\mathbf{L}) = g(\mathbf{L})f(\mathbf{N})\) , and the first assertion follows from the remarks at the beginning of this sub-section.

The group S defined in Proposition 27 is said to be the external semi-direct product of N and L (relative to \(\sigma\) ); we shall generally identify N (resp. L) with the normal subgroup \(j_{1}(N)\) {[}resp. the subgroup \(j_{2}(L)\) {]} of S. If \(\sigma_{y}\) is the identity element of \(\Gamma\) for all \(y \in L\) , we retrieve the usual notion of the product of two groups.

Now let G be a group, and let L and N be two subgroups of G such that LN = NL and such that N is normal in NL, so that for each \(y \in L\) , \(\sigma_{y} : x \to yxy^{-1}\) is an automorphism of N, and \(y \to \sigma_{y}\) is a homomorphism of L into the automorphism group \(\Gamma\) of N. It follows then from Proposition 27 that if S is the external semi-direct product of N and L (relative to \(\sigma\) ), then \(h : (x, y) \to xy\) is a homomorphism of S into G. h is bijective if and only if we have \(N \cap L = \{e\}\) and NL = G (Corollary to Proposition 27); G is then said to be the semi-direct product of its normal subgroup N and its subgroup L, and we often identify G with S by means of h.

PROPOSITION 28. Let L, N be two topological groups, let \(y \to \sigma_y\) be a homomorphism of L into the group of automorphisms \(\Gamma\) of the (non-topological) group structure of N; and suppose that the mapping \((x, y) \to \sigma_y(x)\) of \(N \times L\) into N is continuous. Then:

\begin{enumerate}
\def\labelenumi{\arabic{enumi})}
\item
  On the external semi-direct product S of N and L, relative to \(\sigma\) , the product of the topologies of N and L is compatible with the group structure; the canonical injections \(j_{1}: \mathrm{N} \to \mathrm{S}\) and \(j_{2}: \mathrm{L} \to \mathrm{S}\) are isomorphisms of the topological groups N and L respectively onto the subgroups \(j_{1}(\mathrm{N})\) and \(j_{2}(\mathrm{L})\) of S, and \(pr_{2}\) is a strict morphism of S onto L.
\item
  Let \(f: \mathbf{N} \to \mathbf{G}\) and \(g: \mathbf{L} \to \mathbf{G}\) be two continuous homomorphisms into a topological group \(\mathbf{G}\) , satisfying (4); then the homomorphism
\end{enumerate}

\[
(x, y) \rightarrow f (x) g (y)
\]

of S into G is continuous.

This is an immediate consequence of the definitions and of the properties of the product topology.

The topological group S thus defined is said to be the external topological semi-direct product of N and L (relative to \(\sigma\) ); notice that the condition imposed on \(\sigma\) implies that L operates continuously on the left on N according to the external law \((x, y) \to \sigma_{y}(x)\) {[}no. 4{]}.

Now let G be a topological group and let N and L be two subgroups of G such that G is the semi-direct product of N and L, qua non-topological group; it is then clear that the mapping \((x, y) \to \sigma_{y}(x)\) is continuous on \(N \times L\) , and that the canonical bijective homomorphism

\[
h: (x, y) \rightarrow x y
\]

of S onto G is continuous. But this homomorphism is not necessarily bicontinuous; when it is bicontinuous, G is said to be the topological semi-direct product of N and L. For this to be so it is necessary and sufficient that, if \(p: G \to N\) and \(q: G \to L\) are the mappings which make correspond to \(z \in G\) the unique elements \(p(z) \in N\) and \(q(z) \in L\) such that \(z = p(z)q(z)\) , then one of the mappings p, q is continuous (in which case both are continuous). It comes to the same thing to say that the restriction to L of the canonical mapping \(G \to G/N\) is an isomorphism of the topological group L onto the topological group G/N.

